\documentclass[reqno]{amsart}
\usepackage[T1]{fontenc}
\usepackage{amsmath,amssymb,amsthm}
\usepackage{xr-hyper}
\usepackage[hidelinks]{hyperref}
\externaldocument[ps-]{20261125T010500--poincare-series-and-near-orthogonality__workbook}
\externaldocument[da-]{20261126T030300--the-data__workbook}
\externaldocument[ob-]{20261126T040100--spaces-forms-and-twists__workbook}
\externaldocument[bp-]{20261126T020210--bessel-poincare-series-and-their-gram-operators__workbook}
\externaldocument[m1-]{20261126T051900--unfolding-twice-gram-operators-as-sums-over-rational-points__workbook}
\externaldocument[bd-]{20261126T081400--black-box-discrete-series-and-the-comparison-of-traces__workbook}

\newcounter{chap}
% Chapter numbers, generated from 20261125T012000--grh-workbook__guide.org
% by scripts/make-chapters.py.  Do not edit: change the order of the
% index and rerun the script (./preview-site.sh and CI do this).
%
% Usage in a sheet, after \newcounter{chap}:
%   % Chapter numbers, generated from 20261125T012000--grh-workbook__guide.org
% by scripts/make-chapters.py.  Do not edit: change the order of the
% index and rerun the script (./preview-site.sh and CI do this).
%
% Usage in a sheet, after \newcounter{chap}:
%   % Chapter numbers, generated from 20261125T012000--grh-workbook__guide.org
% by scripts/make-chapters.py.  Do not edit: change the order of the
% index and rerun the script (./preview-site.sh and CI do this).
%
% Usage in a sheet, after \newcounter{chap}:
%   \input{chapters}\chapnumber{<sheet ID>}
\makeatletter
\providecommand\chapters@set[2]{\@namedef{chapters@#1}{#2}}
\chapters@set{20261125T001000}{1}% zeros-to-ramp-sums
\chapters@set{20261125T001100}{2}% inert-primes-and-estimate-b
\chapters@set{20261125T005000}{3}% galois-fourier-and-averaging
\chapters@set{20261125T005500}{4}% columns-from-a-sphere
\chapters@set{20261125T000000}{5}% slope-method
\chapters@set{20261126T030100}{6}% a-twist-for-each-prime
\chapters@set{20261125T003000}{7}% the-chain-and-the-square-root
\chapters@set{20261125T010000}{8}% rank-lemma-and-zero-estimate
\chapters@set{20261125T010500}{9}% poincare-series-and-near-orthogonality
\chapters@set{20261124T223000}{10}% elliptic-curves-for-the-proof
\chapters@set{20261125T009500}{11}% fock-model-and-theta-intertwiners
\chapters@set{20261124T223100}{12}% complex-multiplication-and-inert-reduction
\chapters@set{20261125T008000}{13}% one-cm-determinant
\chapters@set{20261124T214500}{14}% supersingular-theta-divisibility
\chapters@set{20261124T214600}{15}% the-isotropic-plane
\chapters@set{20261125T013700}{16}% pairings-on-a-formal-group-of-height-two
\chapters@set{20261125T013500}{17}% hecke-nondegeneracy-on-the-tree
\chapters@set{20261126T030200}{18}% the-bound-is-attained
\chapters@set{20261125T009000}{19}% kronecker-limit-and-chowla-selberg
\chapters@set{20261125T002000}{20}% definite-theta-lifts
\chapters@set{20261125T015500}{21}% siegel-weil-in-the-smallest-case
\chapters@set{20261125T002100}{22}% definite-theta-lifts-and-norms
\chapters@set{20261125T016500}{23}% hecke-operators-and-theta-lifts
\chapters@set{20261125T016700}{24}% divisibility-along-the-chain
\chapters@set{20261125T004000}{25}% denominators-and-sha
\chapters@set{20261125T004100}{26}% tate-twists-and-sha
\chapters@set{20261126T020430}{27}% the-minus-part-by-gauss-sums
\chapters@set{20261125T014100}{28}% congruence-modules
\chapters@set{20261125T014500}{29}% eisenstein-congruences-and-selmer-classes
\chapters@set{20261125T014700}{30}% pseudocharacters-and-the-length-bound
\chapters@set{20261125T011000}{31}% galois-representations-black-box
\chapters@set{20261126T030300}{32}% the-data
\chapters@set{20261126T020210}{33}% bessel-poincare-series-and-their-gram-operators
\chapters@set{20261126T030400}{34}% the-seesaw-identity
\chapters@set{20261126T030500}{35}% divisibility-of-the-composite-lift
\chapters@set{20261126T030600}{36}% theorem-a-short-forms
\chapters@set{20261126T030700}{37}% the-determinant-of-a-chain
\chapters@set{20261126T030800}{38}% the-cost-of-one-lift
\chapters@set{20261126T030900}{39}% denominators-of-projections-to-lifts
\chapters@set{20261126T031000}{40}% twists-with-few-congruences
\chapters@set{20261126T031100}{41}% theorem-c-and-the-main-inequality
\providecommand\chapnumber[1]{%
  \@ifundefined{chapters@#1}%
    {\PackageError{chapters}{Sheet #1 is not listed in the index}%
      {List it in 20261125T012000--grh-workbook__guide.org\MessageBreak
       and run scripts/make-chapters.py.}}%
    {\edef\thechap{\@nameuse{chapters@#1}}}}
\makeatother
\chapnumber{<sheet ID>}
\makeatletter
\providecommand\chapters@set[2]{\@namedef{chapters@#1}{#2}}
\chapters@set{20261125T001000}{1}% zeros-to-ramp-sums
\chapters@set{20261125T001100}{2}% inert-primes-and-estimate-b
\chapters@set{20261125T005000}{3}% galois-fourier-and-averaging
\chapters@set{20261125T005500}{4}% columns-from-a-sphere
\chapters@set{20261125T000000}{5}% slope-method
\chapters@set{20261126T030100}{6}% a-twist-for-each-prime
\chapters@set{20261125T003000}{7}% the-chain-and-the-square-root
\chapters@set{20261125T010000}{8}% rank-lemma-and-zero-estimate
\chapters@set{20261125T010500}{9}% poincare-series-and-near-orthogonality
\chapters@set{20261124T223000}{10}% elliptic-curves-for-the-proof
\chapters@set{20261125T009500}{11}% fock-model-and-theta-intertwiners
\chapters@set{20261124T223100}{12}% complex-multiplication-and-inert-reduction
\chapters@set{20261125T008000}{13}% one-cm-determinant
\chapters@set{20261124T214500}{14}% supersingular-theta-divisibility
\chapters@set{20261124T214600}{15}% the-isotropic-plane
\chapters@set{20261125T013700}{16}% pairings-on-a-formal-group-of-height-two
\chapters@set{20261125T013500}{17}% hecke-nondegeneracy-on-the-tree
\chapters@set{20261126T030200}{18}% the-bound-is-attained
\chapters@set{20261125T009000}{19}% kronecker-limit-and-chowla-selberg
\chapters@set{20261125T002000}{20}% definite-theta-lifts
\chapters@set{20261125T015500}{21}% siegel-weil-in-the-smallest-case
\chapters@set{20261125T002100}{22}% definite-theta-lifts-and-norms
\chapters@set{20261125T016500}{23}% hecke-operators-and-theta-lifts
\chapters@set{20261125T016700}{24}% divisibility-along-the-chain
\chapters@set{20261125T004000}{25}% denominators-and-sha
\chapters@set{20261125T004100}{26}% tate-twists-and-sha
\chapters@set{20261126T020430}{27}% the-minus-part-by-gauss-sums
\chapters@set{20261125T014100}{28}% congruence-modules
\chapters@set{20261125T014500}{29}% eisenstein-congruences-and-selmer-classes
\chapters@set{20261125T014700}{30}% pseudocharacters-and-the-length-bound
\chapters@set{20261125T011000}{31}% galois-representations-black-box
\chapters@set{20261126T030300}{32}% the-data
\chapters@set{20261126T020210}{33}% bessel-poincare-series-and-their-gram-operators
\chapters@set{20261126T030400}{34}% the-seesaw-identity
\chapters@set{20261126T030500}{35}% divisibility-of-the-composite-lift
\chapters@set{20261126T030600}{36}% theorem-a-short-forms
\chapters@set{20261126T030700}{37}% the-determinant-of-a-chain
\chapters@set{20261126T030800}{38}% the-cost-of-one-lift
\chapters@set{20261126T030900}{39}% denominators-of-projections-to-lifts
\chapters@set{20261126T031000}{40}% twists-with-few-congruences
\chapters@set{20261126T031100}{41}% theorem-c-and-the-main-inequality
\providecommand\chapnumber[1]{%
  \@ifundefined{chapters@#1}%
    {\PackageError{chapters}{Sheet #1 is not listed in the index}%
      {List it in 20261125T012000--grh-workbook__guide.org\MessageBreak
       and run scripts/make-chapters.py.}}%
    {\edef\thechap{\@nameuse{chapters@#1}}}}
\makeatother
\chapnumber{<sheet ID>}
\makeatletter
\providecommand\chapters@set[2]{\@namedef{chapters@#1}{#2}}
\chapters@set{20261125T001000}{1}% zeros-to-ramp-sums
\chapters@set{20261125T001100}{2}% inert-primes-and-estimate-b
\chapters@set{20261125T005000}{3}% galois-fourier-and-averaging
\chapters@set{20261125T005500}{4}% columns-from-a-sphere
\chapters@set{20261125T000000}{5}% slope-method
\chapters@set{20261126T030100}{6}% a-twist-for-each-prime
\chapters@set{20261125T003000}{7}% the-chain-and-the-square-root
\chapters@set{20261125T010000}{8}% rank-lemma-and-zero-estimate
\chapters@set{20261125T010500}{9}% poincare-series-and-near-orthogonality
\chapters@set{20261124T223000}{10}% elliptic-curves-for-the-proof
\chapters@set{20261125T009500}{11}% fock-model-and-theta-intertwiners
\chapters@set{20261124T223100}{12}% complex-multiplication-and-inert-reduction
\chapters@set{20261125T008000}{13}% one-cm-determinant
\chapters@set{20261124T214500}{14}% supersingular-theta-divisibility
\chapters@set{20261124T214600}{15}% the-isotropic-plane
\chapters@set{20261125T013700}{16}% pairings-on-a-formal-group-of-height-two
\chapters@set{20261125T013500}{17}% hecke-nondegeneracy-on-the-tree
\chapters@set{20261126T030200}{18}% the-bound-is-attained
\chapters@set{20261125T009000}{19}% kronecker-limit-and-chowla-selberg
\chapters@set{20261125T002000}{20}% definite-theta-lifts
\chapters@set{20261125T015500}{21}% siegel-weil-in-the-smallest-case
\chapters@set{20261125T002100}{22}% definite-theta-lifts-and-norms
\chapters@set{20261125T016500}{23}% hecke-operators-and-theta-lifts
\chapters@set{20261125T016700}{24}% divisibility-along-the-chain
\chapters@set{20261125T004000}{25}% denominators-and-sha
\chapters@set{20261125T004100}{26}% tate-twists-and-sha
\chapters@set{20261126T020430}{27}% the-minus-part-by-gauss-sums
\chapters@set{20261125T014100}{28}% congruence-modules
\chapters@set{20261125T014500}{29}% eisenstein-congruences-and-selmer-classes
\chapters@set{20261125T014700}{30}% pseudocharacters-and-the-length-bound
\chapters@set{20261125T011000}{31}% galois-representations-black-box
\chapters@set{20261126T030300}{32}% the-data
\chapters@set{20261126T020210}{33}% bessel-poincare-series-and-their-gram-operators
\chapters@set{20261126T030400}{34}% the-seesaw-identity
\chapters@set{20261126T030500}{35}% divisibility-of-the-composite-lift
\chapters@set{20261126T030600}{36}% theorem-a-short-forms
\chapters@set{20261126T030700}{37}% the-determinant-of-a-chain
\chapters@set{20261126T030800}{38}% the-cost-of-one-lift
\chapters@set{20261126T030900}{39}% denominators-of-projections-to-lifts
\chapters@set{20261126T031000}{40}% twists-with-few-congruences
\chapters@set{20261126T031100}{41}% theorem-c-and-the-main-inequality
\providecommand\chapnumber[1]{%
  \@ifundefined{chapters@#1}%
    {\PackageError{chapters}{Sheet #1 is not listed in the index}%
      {List it in 20261125T012000--grh-workbook__guide.org\MessageBreak
       and run scripts/make-chapters.py.}}%
    {\edef\thechap{\@nameuse{chapters@#1}}}}
\makeatother

\ifcsname chapters@20261126T081500\endcsname
  \chapnumber{20261126T081500}
\else
  \renewcommand{\thechap}{?}
\fi
\renewcommand{\thesection}{\thechap.\arabic{section}}

\theoremstyle{definition}
\newtheorem{problem}{Problem}[section]
\newtheorem{definition}{Definition}[section]
\theoremstyle{plain}
\newtheorem{theorem}{Theorem}[section]
\theoremstyle{remark}
\newtheorem*{leanremark}{In the Lean formalization}

\newcommand{\C}{\mathbb C}
\newcommand{\R}{\mathbb R}
\newcommand{\ii}{\mathrm i}
\newcommand{\Uc}{\mathcal U}
\newcommand{\cf}{\mathsf c}
\newcommand{\Per}{\mathrm{Per}}
\newcommand{\diag}{\operatorname{diag}}
\newcommand{\tr}{\operatorname{tr}}
\newcommand{\vol}{\operatorname{vol}}
\newcommand{\tl}[1]{\texttt{#1}}
\setlength{\emergencystretch}{3em}

\title{A Bessel function at a real place}
\author{}
\date{November 26, 2026}

\begin{document}
\maketitle
\refstepcounter{chap}\label{chap}

The real Bessel functions in Theorem~\ref{m1-thm:bessel-system} supply the local factors in the Gram operators for Theorem~\ref{bp-thm:B}.  The proof averages a matrix coefficient over the Bessel subgroup, using a position at which the average is large enough to normalize.  The half-plane and two unitary groups give models for choosing the position.  The period determines the squared norm and the correlation of the average.

\paragraph{Prerequisites.}
Use Chapter~\ref{ps-chap} for the half-plane and Chapters~\ref{da-chap}, \ref{ob-chap} and~\ref{bp-chap} for the data, weights and Bessel subgroup, including the \emph{Bessel tuple} \((v_1,\dots,v_\varrho,z)\) of Section~\ref{bp-sec:subgroup}.  Assume the discrete-series facts of Chapter~\ref{bd-chap} and Schur's lemma \cite{Se}.  Inner products are linear in the first variable, as in Chapter~\ref{m1-chap}.

\section{The half-plane}\label{sec:half-plane}

The first model recovers the seed of a holomorphic Poincar\'e series from a matrix coefficient.  Use integers \(k\ge2\) and \(m\ge1\) for the weight and index in this section.  On the upper half-plane let \(e(z)=\exp(2\pi\ii z)\) and \(j(g,z)=g_{21}z+g_{22}\) for \(g\in\mathrm{SL}_2(\R)\) acting by fractional linear transformations.  Write
\[
 n(w)=\begin{pmatrix}1&w\\0&1\end{pmatrix},\qquad
 a(y)=\diag(\sqrt y,1/\sqrt y),\qquad
 N_{\R}=\{n(w):w\in\R\}.
\]
The representation \((\pi_k(g)f)(z)=j(g^{-1},z)^{-k}f(g^{-1}z)\) and its unit kernel vector \(v_{\ii}\) are as in Chapter~\ref{ps-chap}, with the slash action on \(\mathrm{SL}_2(\R)\).  Use the conjugate coefficient \(\cf(g)=\langle v_{\ii},\pi_k(g)v_{\ii}\rangle\) and formal degree \(d=(k-1)/(4\pi)\) for \(dw\,dy/y^2\) times stabilizer probability (Problem~\ref{ps-pr:formal-degree}).

\begin{problem}\label{pr:half-coeff}
For \(w\in\R\) and \(y>0\),
\[
 \cf(n(w)a(y))=y^{k/2}\bigl((1+y-\ii w)/2\bigr)^{-k}.
\]
Also \(\cf(1)=1\) and \(\cf(g^{-1})=\overline{\cf(g)}\).
\end{problem}

Fix \(y_0>0\) and use \(y_*=a(y_0)\) as the position.  Average against the inverse of the character \(n(w)\mapsto e(mw)\) of \(N_{\R}\), writing \(n(-w)\) as the integration variable.
\[
 W_{y_*}(g)=\int_{\R}\cf(y_*^{-1}n(-w)g)e(mw)\,dw .
\]

\begin{problem}\label{pr:half-integral}
At \(g=a(y)\) the average is
\[
 \frac{(4\pi)^k(yy_0)^{k/2}m^{k-1}}{(k-1)!}
       e^{-2\pi m(y+y_0)} .
\]
Check by machine that at \((k,m,y_0,y)=(4,1,0.3,0.5)\) the value is \(0.613576\) to six decimal places.
\end{problem}

\begin{problem}\label{pr:seed}
For arbitrary \(g\), the average is \(e(m\,g\ii)j(g,\ii)^{-k}\) times \((4\pi)^k y_0^{k/2}m^{k-1}e^{-2\pi my_0}/(k-1)!\).
\end{problem}

Write \(P(s)=W_{y_*}(y_*)/y_0\), where \(s=4\pi my_0\).  The next problem relates the period to the squared norm on \(N_{\R}\backslash\mathrm{SL}_2(\R)\).

\begin{problem}\label{pr:half-norm}
One has \(d\int_0^\infty |W_{y_*}(a(y))|^2\,dy/y^2=W_{y_*}(y_*)\), with common value about \(0.776107\) when \((k,m,y_0)=(4,1,0.3)\).  The set \(\{w: |w|/y_0\le1\}\) has volume \(2y_0\), and
\[
 P(s)=\frac{4\pi s^{k-1}e^{-s}}{(k-1)!}>0.
\]
At \(s=1\) check \(P(1)\approx1.16\cdot10^{-7}\) for \(k=12\) and \(4.95\cdot10^{-156}\) for \(k=100\).
\end{problem}

The next problems estimate the period at a position where the phase is stationary.

\begin{problem}\label{pr:stationary}
At \(g=y_*\), the integrand in the coordinate \(u=w/(2y_0)\) is \(Z(u)=(1+\ii u)^{-k}e^{\ii s u}\).  Its continuous argument vanishing at \(u=0\) has derivative zero there exactly when \(s=k\).  At that position,
\[
 |Z(u)|=(1+u^2)^{-k/2},\qquad |\arg Z(u)|\le k|u|^3/3 .
\]
\end{problem}

\begin{problem}\label{pr:stirling}
At \(s=k\) one has \(P(k)=4\pi k^k e^{-k}/k!\) and \(2y_0\le\sqrt{k}\,W_{y_*}(y_*)\).  Check that \(P(k)/\sqrt{8\pi/k}\) is approximately \(0.9595,0.9794,0.99308,0.999167\) for \(k=2,4,12,100\), respectively.
\end{problem}

\begin{problem}\label{pr:cone}
At \(s=k\), write \(d_k=k^{-1/3}\).  On \(|u|\le d_k\) one has \(|Z(u)|\le2\operatorname{Re}Z(u)\), and \(Z(-u)=\overline{Z(u)}\).  Deduce
\[
 \int_{\R} Z(u)\,du\ \ge\
 \tfrac12\int_{|u|\le d_k}|Z(u)|\,du-\int_{|u|>d_k}|Z(u)|\,du .
\]
Check both sides by machine.  For \(k=12\) they are about \(0.71859\) and \(0.18324\), and for \(k=100\) they are \(0.25045\) and \(0.11321\).  For \(k=4\) the right side is about \(-0.0511\).
\end{problem}

\begin{problem}\label{pr:tail}
The coefficient has \(|\cf(a(e^{2t}))|=\cosh(t)^{-k}\), about \(0.60703\) at \(k=100,t=0.1\), while \(e^{-kt^2/2}\approx0.60653\).  The measure \(|W_{y_*}(a(y))|^2\,dy/y^2\) has density proportional to \(y^{k-2}e^{-4\pi my}\).  At \(s=k\) and \(k=12\), the shares of its mass on \(y>4y_0\) and \(y<y_0/4\) are about \(3.2\cdot10^{-11}\) and \(2.9\cdot10^{-4}\).
\end{problem}

\section{A circle and a flag}\label{sec:positions}

On the half-plane changing a height matches the phase derivatives at the identity.  For \(G=U(1,1)\) preserving \(J=\diag(1,-1)\), the change is a conjugation of the circle \(R=\{\diag(w,1):|w|=1\}\).  Use probability measure \(dw\) on the circle.  In this model \(\mathsf a\ge1\) and \(\mathsf N\ge2\) are integers, and the coefficient formula of Theorem~\ref{bd-thm:fj} gives \(\cf(g)=\det(g)^{-\mathsf a}g_{22}^{-\mathsf N}\) for lowest \(K\)-type \((-\mathsf a,-\mathsf a-\mathsf N)\).  Pair the coefficient against \(w^{t_0}\) with \(t_0\) a nonnegative integer.  For columns write \(I^\dagger=I^*J\).

\begin{problem}\label{pr:circle-period}
The period at the identity vanishes.  For a real \(\mathsf b\ge0\), let \(I_0=(\sqrt{1+\mathsf b},\sqrt{\mathsf b})^T\) and use the inverse of the real boost with first column \(I_0\) as the position.  The conjugated circle is \(g(w)=1+(w-1)I_0I_0^\dagger\), with \(\det g(w)=w\) and \(g(w)_{22}=1+\mathsf b-\mathsf b w\).  For \(j=\mathsf a+t_0\),
\[
 \int_{|w|=1}\cf(g(w))w^{-t_0}\,dw
 =\binom{\mathsf N+j-1}{j}\frac{\mathsf b^j}{(1+\mathsf b)^{\mathsf N+j}} .
\]
For \((\mathsf a,\mathsf N,t_0)=(2,3,0)\) the periods at \(\mathsf b=1/3,2/3,1\) are \(81/512,648/3125,3/16\); for \((2,3,1)\) the last two are \(432/3125,5/32\).
\end{problem}

\begin{problem}\label{pr:circle}
The period is largest at \(\mathsf b=j/\mathsf N\), equivalently when
\(I_0^\dagger D I_0=t_0\) for \(D=\diag(-\mathsf a,-\mathsf a-\mathsf N)\).  The same equation says that the argument of \(\cf(g(e^{\ii\theta}))e^{-\ii t_0\theta}\) is stationary at \(\theta=0\).  As \(j,\mathsf N\) tend to infinity in fixed ratio the maximum is asymptotic to \((2\pi j(j+\mathsf N)/\mathsf N)^{-1/2}\).  For \((j,\mathsf N)=(2,3),(20,30),(200,300)\) the ratios to the asymptotic value are \(0.94897,0.99474,0.99947\) to five decimal places.
\end{problem}

The compression \(I_0^\dagger D I_0\) is the weight seen on the embedded line.  For the role of such weight restrictions in the orbit method, see Nelson and Venkatesh \cite{NV}.

In \(U(2,2)\) the position also involves an isotropic line.  Use \(J=\diag(1,-1,1,-1)\) and the same adjoint notation.  In this model let \((\mu_i)=(-1,-2,-3,-4)\), \(p(z)=\prod_i(z-\mu_i)\), \(q(z)=z\) (the smaller weight is \(0\)), and \(D=\diag(\mu_i)\).  For a polynomial \(F\) write \([F]_i=F(\mu_i)|p'(\mu_i)q(\mu_i)|^{-1/2}\).

\begin{problem}\label{pr:flag}
One has \((p'(\mu_i))=(6,-2,2,-6)\).  For \(v_0=[z]\), \(v_1=[z^2]\) and \(I_0=\sqrt{24}[1]\), the matrix of pairings in that order is \(\diag(0,-1,1)\), and
\[
 Dv_0=v_1,\qquad DI_0=\sqrt{24}v_0,\qquad I_0^\dagger DI_0=0 .
\]
\end{problem}

\begin{problem}\label{pr:flag-phase}
Let \(Y=-\ii(v_0v_1^\dagger+v_1v_0^\dagger)\).  Then \(Y^3=0\) and \(g(t)=1+tY+t^2Y^2/2\), \(t\in\R\), is a one-parameter subgroup of \(U(2,2)\) fixing \(v_0\) and \(I_0\).  It has \(g(t)v_1=v_1+\ii t v_0\), \(\tr(g(t)-1)=0\), and \(\tr(D(g(t)-1))=2\ii t\), giving \(0.2\ii\) at \(t=0.1\).
\end{problem}

\section{The period bound for the data}\label{sec:period}

The general period bound gives a position at a real place of the data.  The model calculations illustrate phase cancellation at small weights.  The bound for the data uses much larger gaps.

Fix a place \(\tau\).  Write \(G=G_\tau=U(Y_0)_\tau\), \(H=U(X_0)_\tau\), and \(R=\Uc_\tau\rtimes H\) as in Section~\ref{bp-sec:subgroup}.  The \emph{sign basis} is obtained from the rescaled basis by a Euclidean-unitary change that diagonalizes the form in the order of the signs for the parameter \(2\varsigma_\tau\).  Use \(ds=du\,dh\) with \(dh\) Haar probability.  In the rescaled basis, \(du\) is Lebesgue measure on the entries in row \(v_a\) and columns \(v_b,v'_b\) for \(b>a\), and on the entries in that row in the \(X_0\) columns and column \(z\), together with the imaginary part of the entry in column \(v'_a\).  The measure \(du\) is used in the covolume of Definition~\ref{bp-def:scaled}.  Use \(dg\) as in Theorem~\ref{bd-thm:degree} and \(dy\) on \(\Uc_\tau\backslash G\) determined by \(du,dg\).

The character from Section~\ref{bp-sec:poincare} is \(\exp(\sigma_R-\overline{\sigma_R})\) at \(\tau\), where \(\sigma_R\) is the sum of the entries just above the diagonal on the rescaled Bessel tuple.  With the chosen complex embedding at \(\tau\), the \(j\)-th \emph{member of the Bessel tuple} is multiplied by \(\bigl(M\tau(\delta_E)/(2\pi\ii)\bigr)^{\varrho+1-j}/\sqrt N\).  The character therefore uses \(\sigma_R=2\pi\ii\lambda(u)/(M\tau(\delta_E))\), with \(\lambda(u)\) the lattice-basis sum of Section~\ref{bp-sec:subgroup} evaluated at the real place.

Let \(\pi_l\) have doubled Harish-Chandra parameter \(l=2\varsigma_\tau\) in the sign basis.  For a unit highest-weight vector \(u_\tau\) in the lowest \(K\)-type, take \(\cf_\tau(y)=\langle\pi_l(y)u_\tau,u_\tau\rangle\) in that basis, with formal degree \(d_\tau\) (Theorems~\ref{bd-thm:ds} and~\ref{bd-thm:degree}).  Take the irreducible unitary representation \(\sigma_\tau\) of \(H\) on \(\C^{I_\tau}\) with the source highest weight given by the diagram of Definition~\ref{ob-def:diagram}, and let \(v\) be a unit highest-weight vector.  Write \(\pi(uh)=\psi_{\Uc}(u)\sigma_\tau(h)\).

\begin{definition}\label{def:average}
For a position \(y_*\in G\), the \emph{Bessel average} at \(v\) is
\[
 W_{y_*,v}(y)=\int_R\cf_\tau(y_*^{-1}sy)\,\pi(s)^{-1}v\,ds .
\]
The integral converges absolutely for the data, uniformly in the sense that the integral of the norm is bounded as \(y\) varies.  This convergence estimate is read from the proof.
\end{definition}

\begin{definition}\label{def:period}
The period at \(y_*\) is \(\Per(y_*)=\langle W_{y_*,v}(y_*),v\rangle\), the integral of \(\cf_\tau(y_*^{-1}sy_*)\overline{\langle\pi(s)v,v\rangle}\) over \(R\).  Write
\[
 V_1(y_*)=\vol_R\{s:\|y_*^{-1}sy_*-1\|\le1\},
\]
using the operator norm in the sign basis.  The \emph{Bessel length} of \(y\) is \(\inf_{s\in R}(\|sy\|+1)\), with the operator norm in the rescaled basis.
\end{definition}

\begin{theorem}[Lower bound for the period]\label{thm:period}
For the data, at each \(\tau\) there is \(y_*\in G_\tau\) of Bessel length at most \(b^{5n}\), \(n=\varrho+1\), with \(V_1(y_*)\le b^{60(r+k)^2}|\Per(y_*)|\).  Here \(r,k\) are the ranks of \(Y_0,X_0\).
\end{theorem}

\noindent\emph{Status}.  Read from the proof; not proved here.  At the position used in the proof, \(g=y_*^{-1}sy_*\), for \(s\in R\), sends each member of the translated Bessel tuple \((y_*^{-1}v_1,\dots,y_*^{-1}v_\varrho,y_*^{-1}z)\) to itself plus a combination of earlier vectors.  Let \(h\) be the compact projection of \(s\).  Let \(D_l\) be the diagonal matrix of the highest weight of the lowest \(K\)-type of \(\pi_l\) in the sign basis, and \(D_u\) that of the highest weight of \(\sigma_\tau\) in the standard basis of \(X_0\).  In these coordinates,
\(\tr(D_l(g-1))-\tr(D_u(h-1))=\sigma_R-\overline{\sigma_R}\), where the sum \(\sigma_R\) is additive under multiplication in \(R\).  Write \(F(s)\) for the period integrand of Definition~\ref{def:period}.  The counterpart of Problem~\ref{pr:cone} is
\[
 |F(s)|\le2\operatorname{Re}F(s)+
  6b^{-150(r+k)^2}(2/\|g\oplus h\|)^{9(r+k)^2},
\]
with the last power integrable on \(R\).  The norm of \(g\oplus h\) is the operator norm on the direct sum.

\section{Norm and correlation}\label{sec:norm}

The following calculations relate the period to the norm and correlation.  Take a locally compact unimodular group \(G\) and a closed unimodular subgroup \(R=\Uc\rtimes H\), with \(\Uc\) closed and unimodular and \(H\) compact.  Use Haar measures \(du,dg\), \(ds=du\,dh\) with \(dh\) of mass one, and the quotient measure \(dy\) on \(\Uc\backslash G\) determined by \(du,dg\).  In the average use a continuous unitary representation \(\pi\) of \(R\) on a finite-dimensional space \(V\), with \(\Uc\) acting by scalars, and a unit vector \(v\in V\).  Replace \(\cf_\tau\) by a normalized unit-vector coefficient \(\cf\) in \(L^1(G)\) of a square-integrable irreducible representation of formal degree \(d_\tau\).  Schur orthogonality, as in Theorem~\ref{bd-thm:schur} for the real unitary case, gives
\[
 d_\tau\int_G \cf(g_0 a g)\overline{\cf(a)}\,da=\cf(g_0)\cf(g).
\]
Assume \(\sup_y\int_R|\cf(y_*^{-1}sy)|\,ds<\infty\).  For the data the proof supplies this bound and \(\cf_\tau\in L^1(G)\), using the coefficient decay of Theorem~\ref{bd-thm:decay}.  Write \(W=W_{y_*,v}\), with norms on \(\Uc\backslash G\).

\begin{problem}\label{pr:equivariance}
For \(s\in R\), \(W(sy)=\pi(s)W(y)\).  Also \(W\) is bounded and in \(L^1\cap L^2\) on the quotient, in the sense that its norm is bounded and belongs to both spaces.
\end{problem}

\begin{problem}\label{pr:reproduce}
One has \(d_\tau\int_G\cf(a^{-1})W(ya)\,da=W(y)\), and \(W\) is continuous.
\end{problem}

\begin{problem}\label{pr:norm-identity}
One has \(d_\tau\|W\|_2^2=\Per(y_*)\).  In particular the period is real and nonnegative, and \(W\ne0\) at the position of Theorem~\ref{thm:period}.
\end{problem}

\begin{problem}\label{pr:coefficient}
The correlation \(g\mapsto\int_{\Uc\backslash G}\langle W(y),W(yg)\rangle\,dy\) equals \(\overline{\cf(g)}\|W\|_2^2\), and its squared \(L^2(G)\) norm is \(\|W\|_2^4/d_\tau\).  If \(H\) acts irreducibly on \(V\), the coordinate products in an orthonormal basis satisfy \(\int_{\Uc\backslash G}W_i\overline{W_j}\,dy=\delta_{ij}\|W\|_2^2/\dim V\).  Check that the half-plane correlation for \(\|W\|_2=1,\ k=4,\ m=1,\ g=a(2)\) has modulus \(64/81\).
\end{problem}

Scale the average to \(\|W_\tau\|_2^2=\#I_\tau\).  The problems give the transformation law, orthonormality and squared correlation of Theorem~\ref{m1-thm:bessel-system}(a)--(c) for \(W_\tau\).  The cut \(W_\tau^\flat\) multiplies \(W_\tau\) by the cutoff in Bessel length which is \(1\) up to half the radius \(b^{6n+2}\), decreases linearly to \(0\) at the radius and is \(0\) thereafter.  The length is left \(R\)-invariant.

We used the discrete-series facts of Chapter~\ref{bd-chap}, absolute convergence, the period bound and the phase and cone estimates.  The cut and its tail bound in Theorem~\ref{m1-thm:bessel-system}(d) and the common decay majorant for the Gram-operator estimates are read from the proof.

\begin{leanremark}
\raggedright The period bound in \cite{GL} is \tl{IsSteep.exists\_stationary\_point}
as used in \tl{Data.exists\_besselSystem}.  The cone bound uses \tl{IsLaplaceDatum.laplace\_method}.
For the data, the average, convergence and reproduction are \tl{DiscreteSeries.besselAverage},
\tl{integrable\_besselIntegrand} and \tl{convolve\_besselAverage}.
The data's norm identity is \tl{formalDegree\_mul\_integral\_norm\_sq\_besselAverage}
with inner-product slots reversed.  The data's squared correlation and normalized coordinate
pairings are \tl{IsBesselFunction.lintegral\_enorm\_sq\_integral\_inner\_translate}
and \tl{BesselSystem.orthonormal}.
\end{leanremark}

\section*{Hints}

\begin{description}
\item[\ref{pr:half-coeff}, \ref{pr:seed}] Kernel covariance and evaluation at \(\ii\).
\item[\ref{pr:half-integral}] Conjugate Problem~\ref{ps-pr:lipschitz}.  Use SageMath quadrature.
\item[\ref{pr:half-norm}, \ref{pr:tail}] The gamma integral, with scale \(4\pi m\).
\item[\ref{pr:stirling}] Stirling's upper bound with remainder.
\item[\ref{pr:cone}] Use SageMath quadrature of \(|Z|\) and the formula for \(P(k)\).
\item[\ref{pr:circle-period}, \ref{pr:circle}] The binomial series and Stirling's formula.
\item[\ref{pr:flag}] Use partial fractions or Lagrange interpolation at the nodes of \(pq\).
\item[\ref{pr:flag-phase}] Use \(Y^\dagger=-Y\) and the pairings in Problem~\ref{pr:flag}.
\item[\ref{pr:equivariance}, \ref{pr:reproduce}] Haar invariance, unfolding and continuity of \(L^1\) translations.
\item[\ref{pr:norm-identity}, \ref{pr:coefficient}] Use unfolding and the Schur identities.
\end{description}

\begin{thebibliography}{NV}
\bibitem[GL]{GL} L.~Sernau, \emph{GRH for Dirichlet \(L\)-functions, in Lean~4}, a shorter rewriting of the Lean proof, text by Claude, \url{https://github.com/lsernau/grh-lean}, commit \texttt{72fead1}, readable edition \texttt{grh/}.
\bibitem[NV]{NV} P.~D.~Nelson and A.~Venkatesh, \emph{The orbit method and analysis of automorphic forms}, Acta Math.\ \textbf{226} (2021), 1--209.
\bibitem[Se]{Se} J.-P.~Serre, \emph{Linear Representations of Finite Groups}, Springer, 1977.
\end{thebibliography}
\end{document}
